\section{Limitations}
\label{sec:limitations}

We discuss the current limitations of \aris{} and directions for future improvement.

\paragraph{Action space coverage.}
The current 24-action discrete action space has limited coverage of certain failure modes.
In particular, \texttt{flat\_low\_contrast} (where the overall scene has low tonal range)
and \texttt{color\_shift} (where object hue diverges from scene) are not always fully resolved
through material and lighting adjustments alone.
Some objects required 40+ evolution rounds without reaching a keep verdict,
suggesting that these failure modes require actions beyond the current space---
for example, scene background replacement or global tone mapping adjustments.

\paragraph{Dataset scale and category coverage.}
ARIS-Rotate contains 50 objects.
While sufficient to demonstrate the pipeline's feasibility and validate downstream utility,
this scale does not provide the breadth required for a general rotation editing model.
Category coverage is also limited; some specialized or atypical object types
may fall outside the distribution covered by this first validation set.

\paragraph{Rotation-only validation.}
We validate \aris{} on a single editing task (rotation).
The pipeline is designed to generalize to other geometry-controlled editing tasks
(lighting changes, material substitution, scale variation),
but such generalization has not been experimentally verified in this report.

\paragraph{VLM reviewer noise.}
The Qwen3.5-35B-A3B reviewer exhibits score variance of approximately 0.006--0.012 per query,
which can produce inconsistent verdicts for borderline cases.
Repeated review passes or multi-view aggregation can mitigate this,
but reviewer reliability remains a variable in pipeline quality.

\paragraph{No public release.}
The current implementation relies on a specific Blender scene file and server environment
and is not publicly released with this technical report.
A public release is planned upon sufficient generalization of the scene-aware rendering components.

\begin{figure}[t]
  \centering
  % TODO: Insert failure cases figure
  \fbox{\parbox{0.95\linewidth}{\centering\vspace{2cm}
    \textbf{[Figure 6: Failure Cases]}\\
    2--3 objects where the evolution loop could not reach a keep verdict.\\
    Annotate: (a) flat\_low\_contrast unresolved, (b) persistent color\_shift,
    (c) geometry-limited case.
    \vspace{2cm}}}
  \caption{
    Representative failure cases.
    (a) Flat, low-contrast lighting that the action space cannot fully resolve.
    (b) Persistent color shift due to material-scene mismatch beyond adjustable range.
    (c) Geometric distortion from the 3D reconstruction that propagates to rendering artifacts.
  }
  \label{fig:failures}
\end{figure}
